\documentclass[aspectratio=169]{beamer}
\usetheme{Madrid}
\usecolortheme{seagull}

\usepackage[utf8]{inputenc}
\usepackage[T1]{fontenc}
\usepackage{lmodern}
\usepackage{booktabs}
\usepackage{amsmath}
\usepackage{array}
\usepackage{listings}
\usepackage{xcolor}

\title{LINMA2710 Project 2026}
\subtitle{Matrix Operations: Sequential, OpenMP, MPI, OpenCL}
\author{Your Name}
\institute{UCLouvain}
\date{May 2026}

\lstset{
  basicstyle=\ttfamily\scriptsize,
  keywordstyle=\color{blue},
  commentstyle=\color{gray},
  showstringspaces=false,
  breaklines=true,
  frame=single
}

\begin{document}

\begin{frame}
  \titlepage
\end{frame}

\begin{frame}{Project Goal and Structure}
\textbf{Goal:} implement the same matrix API in 4 computing paradigms.

\vspace{0.3cm}
\begin{tabular}{@{}p{2.7cm}p{5.1cm}p{4.2cm}@{}}
\toprule
Class & Data location & Compute model \\
\midrule
\texttt{Matrix} & CPU, contiguous \texttt{std::vector<double>} & Sequential loops + OpenMP pragmas \\
\texttt{DistributedMatrix} & Columns split across MPI ranks & Local compute + collective comms \\
\texttt{MatrixCL} & OpenCL device buffer (GPU/CPU device) & Kernels launched with NDRange \\
\bottomrule
\end{tabular}

\vspace{0.35cm}
\textbf{Main challenge:} keep one logical API while changing where data lives and how operations are executed.
\end{frame}

\begin{frame}{Part 1: Matrix Implementation (CPU)}
\textbf{What was implemented in \texttt{src/matrix.cpp}:}
\begin{itemize}
  \item Constructors, copy, assignment, bounds-checked \texttt{get/set}
  \item \texttt{fill}, \texttt{+}, \texttt{-}, scalar product, matrix product, \texttt{transpose}
  \item \texttt{apply(func)} and in-place \texttt{sub\_mul(scalar, other)}
  \item Dimension checks with exceptions for invalid operations
\end{itemize}

\textbf{Data layout:} row-major contiguous storage
\[
\texttt{data}[i,j] = \texttt{data}[i \cdot \texttt{cols} + j]
\]

\textbf{Complexity:}
\begin{itemize}
  \item Add/sub/scalar/apply/sub\_mul: $O(nm)$
  \item Transpose: $O(nm)$
  \item Matrix multiplication: $O(nmk)$
\end{itemize}
\end{frame}

\begin{frame}{Part 1 Questions: Technical Answers}
\textbf{Copy constructor behavior}
\begin{itemize}
  \item Copy constructor copies \texttt{rows}, \texttt{cols}, and \texttt{std::vector<double> data}
  \item This is a deep copy: modifying the copy does \textbf{not} modify the original
\end{itemize}

\textbf{Why no explicit destructor?}
\begin{itemize}
  \item Rule of zero / RAII: \texttt{std::vector} frees memory automatically
  \item No raw \texttt{new/delete}, so compiler-generated destructor is correct
\end{itemize}

\textbf{Sparse matrix handling}
\begin{itemize}
  \item Current design is dense (good for dense workloads)
  \item For sparse: use CSR/CSC storage + sparse kernels (SpMV/SpMM)
  \item Benefit: memory and compute proportional to nnz instead of $n \times m$
\end{itemize}

\textbf{SIMD speedup}
\begin{itemize}
  \item Possible via compiler auto-vectorization or explicit AVX intrinsics
  \item Typical gain on dense loops: around $1.5\times$ to $4\times$ depending on memory bandwidth and loop structure
\end{itemize}
\end{frame}

\begin{frame}{Part 2: OpenMP Parallelization Strategy}
\textbf{Implementation choices (same \texttt{src/matrix.cpp}):}
\begin{itemize}
  \item OpenMP enabled with \texttt{-fopenmp}
  \item Parallel loops added to:\
  \texttt{fill, +, -, *, transpose, apply, sub\_mul}
  \item Matrix multiplication uses \texttt{\#pragma omp parallel for collapse(2)} on $(i,j)$ loops
\end{itemize}

\textbf{Why this works well:}
\begin{itemize}
  \item Each thread computes independent output elements
  \item No race condition on output matrix entries
  \item Good load balance for medium/large matrices
\end{itemize}

\textbf{Amdahl law reminder:}
\[
S(p)=\frac{1}{(1-f)+\frac{f}{p}}
\]
where $f$ is parallel fraction and $p$ is number of threads.
\end{frame}

\begin{frame}{Part 2 Results: Your OpenMP Benchmark}
\scriptsize
\begin{tabular}{rcccc}
\toprule
Size & $T_1$ (ms) & $T_2$ speedup & $T_4$ speedup & $T_8$ speedup \\
\midrule
32   & 0.04    & 2.11x & 1.57x & 1.75x \\
64   & 0.27    & 2.05x & 3.48x & 3.57x \\
128  & 2.39    & 1.86x & 3.35x & 4.12x \\
256  & 51.52   & 1.99x & 3.94x & 7.64x \\
512  & 575.22  & 2.01x & 4.01x & 5.79x \\
1024 & 9735.27 & 1.99x & 3.99x & 7.99x \\
\bottomrule
\end{tabular}

\vspace{0.2cm}
\textbf{Key observations:}
\begin{itemize}
  \item Near-ideal scaling for large matrices (1024: almost $8\times$ with 8 threads)
  \item 512 case: efficiency drops at 8 threads (measured $72.7\%$)
  \item Small sizes are overhead-dominated (thread startup, scheduling, sync)
\end{itemize}

\textbf{Potential fixes for small sizes:}
\begin{itemize}
  \item Size threshold before parallel region
  \item Tune scheduling/chunk size
  \item Blocked multiplication for better cache reuse
\end{itemize}
\end{frame}

\begin{frame}{Part 3: DistributedMatrix (MPI) Implementation}
\textbf{Partitioning strategy:}
\begin{itemize}
  \item Global matrix split by columns as evenly as possible
  \item Rank $r$ stores local block: columns $[\texttt{startCol},\,\texttt{startCol}+\texttt{localCols})$
\end{itemize}

\textbf{Implemented operations:}
\begin{itemize}
  \item Local-only: \texttt{fill}, \texttt{+}, \texttt{-}, scalar, \texttt{apply}, \texttt{sub\_mul}
  \item Communication-based: \texttt{gather}, \texttt{sum}, \texttt{multiplyTransposed}, \texttt{sync\_matrix}
\end{itemize}

\textbf{multiplyTransposed pattern:}
\begin{enumerate}
  \item Each rank computes local partial matrix product
  \item \texttt{MPI\_Allreduce} with sum to combine partials
\end{enumerate}

\textbf{Performance model:}
\[
T \approx T_{comp} + T_{comm}
\quad\text{with}\quad
T_{comp}\sim O\!\left(\frac{n^3}{p}\right),
\;T_{comm}\sim O(n^2\log p)
\]
\end{frame}

\begin{frame}{Part 3: Communication Profiling and Analysis}
\textbf{Profiling approach used and recommended:}
\begin{itemize}
  \item Measure local computation time per rank for local partial product using high-resolution timers (e.g., std::chrono).
  \item Measure MPI collectives separately: time the \texttt{MPI\_Allreduce} and \texttt{MPI\_Allgatherv} calls using MPI_Wtime.
  \item Report fraction: $\alpha = T_{comm}/(T_{comm}+T_{comp})$.
\end{itemize}

\textbf{Typical results and interpretation:}
\begin{itemize}
  \item For large matrices, $T_{comp}\gg T_{comm}$ so communication overhead is small and speedup approaches ideal.
  \item For moderate sizes or large process counts, communication becomes significant and limits scaling.
  \item Allreduce of an $n\times n$ partial requires sending $O(n^2)$ doubles; cost grows with data size and log(p) factors for algorithms.
\end{itemize}

\textbf{Practical microbenchmark recipe:}
\begin{enumerate}
  \item Fix matrix sizes, vary number of processes $p$.
  \item Time: local partial compute, allreduce, and total end-to-end.
  \item Plot compute vs comm time versus $p$ to identify communication-dominated regime.
\end{enumerate}
\end{frame}

\begin{frame}{Part 3: Expected Speedup and Measured Comparison}
\textbf{Model:} with $p$ processes, if local work scales as $1/p$ and communication is $T_{comm}$, total time:
\[
T(p)=\frac{T_{serial}\cdot f}{p} + T_{serial}\cdot(1-f) + T_{comm}(p)
\]
where $f$ is parallel fraction of compute.

\textbf{Expected speedup:}
\begin{itemize}
  \item Upper bound near $p$ when $T_{comm}\approx 0$ and $f\approx 1$.
  \item More realistic: diminishing returns as $p$ increases — communication and latency limit gains.
\end{itemize}

\textbf{How to compare with measurements:}
\begin{itemize}
  \item Compute measured speedup $S(p)=T(1)/T(p)$.
  \item If $S(p) \ll p$, investigate: increased $T_{comm}$, load imbalance, or memory contention.
\end{itemize}
\end{frame}

\begin{frame}{Part 3: Partitioning vs Gradient-Sync Approach}
\textbf{Current column-split approach (this project):}
\begin{itemize}
  \item Each process stores a contiguous block of columns; local ops operate without communication.
  \item Collective operations (gather, allreduce) needed for global results.
\end{itemize}

\textbf{Gradient-sync (data-parallel) style:}
\begin{itemize}
  \item Each process holds a full model shard or local data; computations are local; gradients aggregated occasionally.
  \item Common in ML: frequent synchronization of smallish gradients or periodic allreduce.
\end{itemize}

\textbf{Trade-offs:}
\begin{itemize}
  \item Column-split: less frequent communication for local elementwise ops, but expensive all-to-all style collective for some linear algebra.
  \item Gradient-sync: lower-latency frequent gradients if they fit network characteristics; can use compression or quantization.
  \item Choice depends on operation sparsity, data size, and communication fabric.
\end{itemize}
\end{frame}

\begin{frame}{Part 3 Questions: Analysis and Comparison}
\textbf{Communication vs computation overhead}
\begin{itemize}
  \item For large matrices, local FLOPs dominate and scaling improves
  \item For smaller matrices or larger $p$, \texttt{MPI\_Allreduce/Allgatherv} dominates runtime
\end{itemize}

\textbf{Expected speedup}
\begin{itemize}
  \item Ideal upper bound: close to $p$ when communication is negligible
  \item Real speedup limited by network latency/bandwidth and collective costs
\end{itemize}

\textbf{Compared to gradient-synchronization style partitioning}
\begin{itemize}
  \item Current project split: model/data already partitioned by columns for matrix algebra
  \item Gradient-sync approach (typical in distributed ML): local forward/backward + periodic global sync
  \item Trade-off: gradient sync has fewer but often larger synchronization phases; this project uses collectives tightly coupled to algebraic operations
\end{itemize}
\end{frame}

\begin{frame}{Part 4: MatrixCL (OpenCL) Implementation}
\textbf{What was implemented in \texttt{src/matrix\_opencl.cpp}:}
\begin{itemize}
  \item Device-side storage with \texttt{cl::Buffer}
  \item Kernel cache with one-time compilation (fill/add/sub\_mul/transpose/matmul)
  \item Host wrappers for API-equivalent operations and copy-back to host
  \item Queue created with \texttt{CL\_QUEUE\_PROFILING\_ENABLE}
\end{itemize}

\textbf{Simple matrix multiplication kernel used:}
\begin{itemize}
  \item One work-item computes one output $C(i,j)$
  \item Inner loop over shared dimension $k$
  \item Correct and readable baseline kernel
\end{itemize}

\textbf{Second, faster version to present (recommended):}
\begin{itemize}
  \item Tiled kernel with local memory (blocked multiplication)
  \item Reduces global memory traffic, improves arithmetic intensity
  \item Best for large matrices and GPUs with enough local memory
\end{itemize}
\end{frame}

\begin{frame}[fragile]{Part 4: Two Kernel Versions — Naive and Tiled}
\textbf{Version 1: naive kernel (already in source)}
\begin{lstlisting}[language=C]
__kernel void matrix_mul(__global const float* A,
                         __global const float* B,
                         __global float* C,
                         int A_rows, int A_cols, int B_cols) {
    int i = get_global_id(0);
    int j = get_global_id(1);
    if (i < A_rows && j < B_cols) {
        float sum = 0.0f;
        for (int k = 0; k < A_cols; ++k) {
            sum += A[i * A_cols + k] * B[k * B_cols + j];
        }
        C[i * B_cols + j] = sum;
    }
}
\end{lstlisting}

\textbf{Version 2: tiled (uses local memory) — example tile size 16}
\begin{lstlisting}[language=C]
#define TS 16
__kernel void matmul_tiled(__global const float* A,
                           __global const float* B,
                           __global float* C,
                           int A_rows, int A_cols, int B_cols) {
    int row = get_global_id(0);
    int col = get_global_id(1);

    __local float As[TS][TS];
    __local float Bs[TS][TS];

    float acc = 0.0f;
    for (int t = 0; t < (A_cols + TS-1)/TS; ++t) {
        int aRow = row;
        int aCol = t*TS + get_local_id(1);
        int bRow = t*TS + get_local_id(0);
        int bCol = col;

        // Load tiles with bounds checks
        As[get_local_id(0)][get_local_id(1)] = (aRow < A_rows && aCol < A_cols) ? A[aRow*A_cols + aCol] : 0.0f;
        Bs[get_local_id(0)][get_local_id(1)] = (bRow < A_cols && bCol < B_cols) ? B[bRow*B_cols + bCol] : 0.0f;

        barrier(CLK_LOCAL_MEM_FENCE);

        for (int k = 0; k < TS; ++k) {
            acc += As[get_local_id(0)][k] * Bs[k][get_local_id(1)];
        }
        barrier(CLK_LOCAL_MEM_FENCE);
    }

    if (row < A_rows && col < B_cols)
        C[row*B_cols + col] = acc;
}
\end{lstlisting}

\textbf{Why tiled is faster:}
\begin{itemize}
  \item Reuses loaded tiles from fast local memory; reduces global memory traffic
  \item Higher arithmetic intensity per memory load
  \item Tile size tuned to device local memory and register pressure
\end{itemize}
\end{frame}

\begin{frame}{Part 4: Profiling and Energy Measurement Plan}
\textbf{Profiling steps:}
\begin{enumerate}
  \item Enable profiling on the command queue (already set in tests): \texttt{CL\_QUEUE\_PROFILING\_ENABLE}
  \item Enqueue kernel, obtain \texttt{cl::Event}, then use \texttt{getProfilingInfo<CL_PROFILING_COMMAND_START/END>()}
  \item Report host-to-device, kernel, device-to-host times separately
\end{enumerate}

\textbf{Energy measurement:}
\begin{itemize}
  \item Use vendor tools (NVIDIA: \texttt{nvidia-smi --query-gpu=power.draw --format=csv}) or CodeCarbon to sample power during runs
  \item Compute energy-to-solution: integrate power over kernel runtime (Joules)
\end{itemize}

\textbf{Evaluation matrix:}
\begin{itemize}
  \item Vary problem size, compare naive vs tiled runtimes and energy
  \item Report speedup and energy reduction (if any)
\end{itemize}
\end{frame}

\begin{frame}[fragile]{Part 4: Two Kernel Versions (for Presentation)}
\textbf{Version 1: naive (implemented)}
\begin{lstlisting}[language=C]
__kernel void matrix_mul(__global const float* A,
                         __global const float* B,
                         __global float* C,
                         int A_rows, int A_cols, int B_cols) {
  int i = get_global_id(0);
  int j = get_global_id(1);
  if (i < A_rows && j < B_cols) {
    float sum = 0.0f;
    for (int k = 0; k < A_cols; ++k)
      sum += A[i * A_cols + k] * B[k * B_cols + j];
    C[i * B_cols + j] = sum;
  }
}
\end{lstlisting}

\textbf{Version 2: optimized (to discuss)}: tiled local-memory kernel (e.g., tile size 16 or 32), fewer global loads, better throughput.
\end{frame}

\begin{frame}{Part 4 Questions: Profiling and Energy}
\textbf{Profiling approach:}
\begin{itemize}
  \item Use event profiling timestamps from OpenCL queue
  \item Compare kernel execution time naive vs tiled
  \item Also measure end-to-end time (host-device transfers included)
\end{itemize}

\textbf{Expected behavior:}
\begin{itemize}
  \item Naive kernel: memory bandwidth bottleneck
  \item Tiled kernel: higher occupancy/reuse, lower time for large matrices
\end{itemize}

\textbf{Power/energy measurement plan:}
\begin{itemize}
  \item Run fixed-size workloads repeatedly
  \item Log GPU power using CodeCarbon or vendor tools
  \item Report energy-to-solution (Joules) and not only runtime
\end{itemize}
\end{frame}

\begin{frame}{5-Minute Talk Plan}
\textbf{Suggested timing:}
\begin{itemize}
  \item 0:00--0:30: Context and architecture
  \item 0:30--1:30: Part 1 design decisions + question answers
  \item 1:30--2:30: OpenMP method + your benchmark highlights
  \item 2:30--3:30: MPI decomposition + communication/computation trade-off
  \item 3:30--4:30: OpenCL design + two-kernel story + profiling strategy
  \item 4:30--5:00: Conclusion and future optimization steps
\end{itemize}

\textbf{Strong conclusion message:}
\begin{itemize}
  \item Same matrix API, 4 execution models
  \item Parallelization gives major speedups for large workloads
  \item Real performance depends on memory hierarchy and communication costs
\end{itemize}
\end{frame}

\begin{frame}{Backup Slide: Extra Questions You Can Answer}
\begin{itemize}
  \item Why use column partitioning in MPI and not row partitioning?
  \item How would blocked CPU multiplication improve cache behavior?
  \item What numerical precision issues can appear (float vs double)?
  \item How to decide between OpenMP, MPI, and OpenCL for a given workload?
  \item What is the impact of PCIe transfer overhead on GPU speedup?
\end{itemize}
\end{frame}

\end{document}
